\section{Evaluation Data and Reward Signal}
\label{sec:data}

All experiments in this paper share a common evaluation framework: prompt sources, a reward signal, and a feature pipeline.
This section defines these foundations; subsequent sections build on them without repeating the details.

\subsection{Prompt Sources}

We draw prompts from two public sources selected for complementary strengths:

\begin{itemize}[nosep,leftmargin=*]
    \item \textbf{LMSYS Chatbot Arena}~\cite{zheng2023judging}: real user prompts submitted to a public LLM comparison platform. These provide ecologically valid, diverse queries spanning creative writing, reasoning, coding, and structured extraction. All evaluation and online-learning data (1,871 prompts across dev and holdout splits) come from this source.
    \item \textbf{RouteLLM battle corpus}~\cite{ong2024routellm}: 80,000 curated pairwise battles between Mixtral-8x7B and GPT-4-Turbo, released as part of the RouteLLM benchmark. We use this corpus primarily to construct cross-distribution warmup priors for the $K{=}2$ evaluation, matching RouteLLM's own training condition for a controlled comparison against its published baseline. The corpus also serves as the training set for the unsupervised PCA feature projection (Section~\ref{sec:feature_pipeline}).
\end{itemize}

\subsection{Reward Signal}
\label{sec:reward_signal}

At routing time, the system observes \emph{bandit feedback}: only the reward of the selected model is revealed.
Two distinct judging protocols produce these rewards, one per evaluation regime:

\paragraph{$K{=}2$ rewards (pairwise, ternary).}
All 1,871 LMSYS evaluation prompts are judged by GPT-4o at temperature~0 (deterministic) using RouteLLM's MT-Bench-style pairwise template, yielding $r \in \{0, 0.5, 1\}$ (loss/tie/win).
The 80K RouteLLM warmup battles ship with pre-existing pairwise outcome labels from the \texttt{routellm/gpt4\_judge\_battles} dataset; we use these labels directly as warmup rewards.
Any mismatch between the warmup and evaluation reward signals is attenuated by the plasticity factor ($\rho{=}0.1$) that down-scales warmup priors and by the Corralling meta-learner, which down-weights the warmup expert when its predictions diverge from observed rewards.

\paragraph{$K{>}2$ rewards (cross-provider panel).}
The 10 models evaluated in this paper (Table~\ref{tab:data_splits}) are judged on LMSYS Arena prompts by a panel of 3--4 LLMs drawn from four providers: GPT-4o (OpenAI), Claude-3.5-Sonnet (Anthropic), Llama-3.1-405B (Meta), and Gemini-2.5-Pro (Google).
To prevent self-evaluation bias, a judge never evaluates models from its own provider family; each model is therefore scored by the 3 remaining judges.
Judges follow a structured chain-of-thought protocol: each produces an explicit reasoning trace before casting a binary pass/fail vote with a confidence score, reducing snap-judgment bias.
The per-prompt reward is the mean vote across the panel, yielding a continuous signal in $[0,1]$ that preserves evaluative granularity (e.g., a 2-of-3 pass maps to $0.67$ rather than being binarised to $1.0$).
Full reward specification, the CoT prompt template, and known limitations are detailed in Appendix~\ref{app:reward_signal}.

\subsection{Data Splits: Two Evaluation Regimes}
\label{sec:data_splits}

Our evaluation operates at two portfolio scales that require different warmup prior sources, reflecting a natural constraint: available offline data rarely covers every model in a deployment portfolio.
Table~\ref{tab:data_splits} summarises both regimes.

\begin{table}[t]
\centering
\caption{Data splits for two evaluation regimes. The holdout set (750 prompts) is identical across regimes; only the model columns read differ.}
\label{tab:data_splits}
\small
\resizebox{\columnwidth}{!}{%
\begin{tabular}{@{}llrll@{}}
\toprule
\textbf{Regime} & \textbf{Split} & \textbf{Prompts} & \textbf{Source} & \textbf{Purpose} \\
\midrule
\multirow{3}{*}{$K{=}2$}
    & Warmup priors  & 80,000 & RouteLLM battles  & LinUCB priors (+ PCA training, shared) \\
    & Online learning & 1,121 & LMSYS Arena (dev)  & Bandit updates \\
    & Holdout         & 750   & LMSYS Arena        & Frozen evaluation \\
\midrule
\multirow{3}{*}{$K{=}5, 10$}
    & Prior training  & 355   & LMSYS Arena (dev) & Warmup priors for all $K$ models \\
    & Online learning & 533   & LMSYS Arena (dev) & Bandit updates \\
    & Holdout         & 750   & LMSYS Arena       & Frozen evaluation \\
\bottomrule
\end{tabular}%
}
\end{table}

\paragraph{$K{=}2$ (RouteLLM comparison).}
We initialise warmup priors from the 80K RouteLLM battles for Mixtral-8x7B and GPT-4-Turbo.
This is a deliberate \emph{cross-distribution} setup: the battles come from a different user population and time period than the evaluation prompts, matching RouteLLM's own training condition and creating the realistic scenario where offline data does not perfectly match deployment traffic.

\paragraph{$K{=}5$ and $K{=}10$ (multi-model scaling).}
The RouteLLM data covers only 2 models, leaving the remaining $K{-}2$ models without priors.
Instead, we construct warmup priors from a held-out portion of our own LMSYS evaluation data, which covers all 10 portfolio models, providing a \emph{same-distribution} prior reflecting the practitioner workflow of seeding a router with a small batch of historical evaluations.
The 888 development prompts with full model coverage are split via stratified sampling into 355 prior-training and 533 online-learning prompts; the holdout set (750 prompts) is shared with the $K{=}2$ evaluation to ground cross-regime comparisons.
Full construction details and data-integrity checks are in Appendix~\ref{sec:warmup_prior_construction}.

\subsection{Feature Pipeline}
\label{sec:feature_pipeline}

Each prompt is encoded by Sentence-BERT (\texttt{all-MiniLM-L6-v2}~\cite{reimers2019sentence}, 384D) and projected to 32 dimensions via unsupervised PCA trained on the 80K RouteLLM battle prompts.
A bias term is appended, yielding a 33-dimensional context vector.
This projection is applied to prompts from both sources: RouteLLM battles for $K{=}2$ warmup priors, and LMSYS Arena prompts for $K{>}2$ warmup priors and all online learning and evaluation.
The PCA is trained once and reused across all portfolio sizes ($K{=}2, 5, 10$), decoupling feature extraction from warmup-prior construction.
Variance analysis and provenance details are in the supplementary material.

With this evaluation framework and feature pipeline established, we now turn to the empirical properties of LLM routing that motivate our system architecture.
